\documentclass[11pt]{article}
\usepackage[margin=1.1in]{geometry}
\usepackage{amsmath,amssymb,amsthm}
\usepackage{booktabs}
\usepackage{enumitem}
\usepackage[T1]{fontenc}
\usepackage{xcolor}
\usepackage{hyperref}
\hypersetup{colorlinks=true, linkcolor=blue!50!black, citecolor=blue!50!black, urlcolor=blue!50!black}

\newcommand{\vx}{\mathbf{x}}
\newcommand{\vu}{\mathbf{u}}
\newcommand{\vlam}{\boldsymbol{\lambda}}
\newcommand{\vbeta}{\boldsymbol{\beta}}
\newcommand{\Lag}{\mathcal{L}}
\newcommand{\Ham}{H}
\newcommand{\file}[1]{\texttt{#1}}
\newcommand{\code}[1]{\texttt{#1}}
\newtheorem{condition}{Condition}

\title{First-Order Optimality Checks in \file{verify\_common}:\\
Mathematics, Code, and Interpretation}
\author{Michael Casey}
\date{July 2026}

\begin{document}
\maketitle

\begin{abstract}
\noindent\sloppy
The \file{orbit\_transfer/verify\_common/} module implements one generic,
automatic-differentiation--based first-order optimality checker
(\file{foc\_check.m}) that runs identically on all four transfer campaigns
(Earth two-body elliptic$\to$GEO, the same transfer in the Earth--Moon CR3BP,
GTO$\to$tulip, GTO$\to$ELFO), plus a fixed-format report printer
(\file{foc\_report.m}) that every production solve driver now emits. This note
explains each check twice over: (a) the mathematics --- the formal necessary
condition being tested and the intuition behind it --- and (b) the code --- how
to run it on each campaign and how to read every line of the printed report.
The checks run in \emph{report-only burn-in}: a FAIL is an advisory finding and
never demotes a certified solution. This revision incorporates a three-way
external review of the layer (2026-07-25) and the three fixes it produced ---
$\varepsilon$-awareness, mesh-normalized switch regularity, and an
endpoint-corrected transversality --- together with the attribution those
fixes made possible: of five outstanding advisory failures, four were
discretization artifacts and one is real. The cross-campaign status of every
check lives in the register \file{OPTIMALITY\_CERTIFICATION.md} at the
\file{orbit\_transfer} root.
\end{abstract}

\tableofcontents

%=====================================================================
\section{Setting and notation}
\label{sec:setting}

\subsection{The continuous problem}

All four campaigns solve minimum-fuel (or minimum-time / minimum-energy)
low-thrust transfer problems of the common form
\begin{equation}
\min_{\vbeta(\cdot),\,s(\cdot)}
\; J \;=\; \int_{0}^{t_f} \ell\bigl(s(t)\bigr)\,dt,
\qquad
\dot\vx \;=\; f_0(\vx,t) \;+\; \frac{T_{\max}}{m}\,s\,B(\vx)\,\vbeta,
\label{eq:ocp}
\end{equation}
with unit thrust direction $\|\vbeta\|=1$, throttle $s\in[0,1]$, mass flow
$\dot m = -(T_{\max}/c)\,s$, fixed endpoints on the orbital state, free final
mass, and (depending on campaign) fixed or free final time. At the fuel
endpoint of the Bertrand--\'Ep\'enoy homotopy \cite{BertrandEpenoy2002} the
running cost is $\ell(s)=s$, linear in the throttle, so Pontryagin's principle
predicts \emph{bang-bang} control.

Pontryagin's Minimum Principle (PMP) \cite{HMG2004,Lawden1963} introduces the
costate $\vlam(t)$ and the Hamiltonian
\begin{equation}
\Ham(\vx,\vu,\vlam)
\;=\; \vlam^{\!\top} \Bigl( f_0 + \tfrac{T_{\max}}{m}\,s\,B\vbeta \Bigr)
\;+\; \lambda_m\,\dot m \;+\; \ell(s),
\label{eq:ham}
\end{equation}
and asserts, along an optimal trajectory: costate dynamics
$\dot\vlam = -\partial\Ham/\partial\vx$, the pointwise \emph{minimum condition}
$\vu^\ast = \arg\min_{\vu\,\in\,\mathcal U}\Ham$, transversality at free
terminal components, and constancy (autonomous case) or a terminal value
condition (free $t_f$) on $\Ham$. Section~\ref{sec:checks} maps each of these
onto a concrete numerical test.

\subsection{The discrete problem the checks actually see}

Each campaign transcribes \eqref{eq:ocp} by trapezoidal collocation on a mesh
$0=\sigma_1<\dots<\sigma_{N+1}=1$ with steps $h_k=\sigma_{k+1}-\sigma_k$: the
decision vector $z$ stacks the node states
$X\in\mathbb R^{n_x\times(N+1)}$ (column-major), then the node controls
$U\in\mathbb R^{n_u\times(N+1)}$, then any extra scalars (the Earth campaign's
longitude span $\Delta L$). The dynamics enter as \emph{defect} equality
constraints, one $n_x$-vector per interval,
\begin{equation}
d_k \;=\; x_{k+1}-x_k-\tfrac{h_k}{2}\,\Xi\,\bigl(F_k+F_{k+1}\bigr) \;=\;0,
\qquad k=1,\dots,N,
\label{eq:defect}
\end{equation}
where $\Xi$ is the campaign's clock factor ($\Delta L$, the Sundman span, or
the \code{cScale} free-final-time slack state of Betts \cite{Betts2010}). The cone $\|\vbeta_k\|^2=1$, the throttle box
$s_k\in[0,1]$, the state boxes, and the boundary conditions complete the
nonlinear program (NLP)
\begin{equation}
\min_z\; f(z) \quad\text{s.t.}\quad g(z)\in[g_L,\,g_U].
\label{eq:nlp}
\end{equation}
IPOPT returns a primal--dual pair $(z^\ast,\vlam_g)$. Everything in this note
is computed \emph{at that pair}, by differentiating with CasADi the exact
$f$ and $g$ that were solved --- never a hand-transcribed copy of the physics.
That is the design principle of the layer: the checks cannot drift out of sync
with the solver, because they differentiate the solver's own model.

\paragraph{Campaign manifests.}
What varies between campaigns is bookkeeping, not structure. A
\emph{manifest} (\file{foc\_manifest.m}) records it:

\begin{center}
\small
\begin{tabular}{@{}llllll@{}}
\toprule
manifest & $n_x$ (state) & $n_u$ & mass row & time row & horizon \\
\midrule
\code{earth\_mee}     & 7 $[P,e_x,e_y,h_x,h_y,m,t]$ & 4 & 6 & 7 & fixed $t_f$ \\
\code{earth\_cr3bp}   & 7 (same; non-autonomous)      & 4 & 6 & 7 & fixed $t_f$ \\
\code{tulip}          & 8 $[\mathbf r,\mathbf v,m,t]$ & 4 & 7 & 8 & fixed $\tau_f$ \\
\code{elfo\_fuel}     & 9 $[\mathbf r,\mathbf v,m,t,c]$ & 4 & 7 & 8 & free $t_f$ (\code{cScale}) \\
\code{elfo\_mintime}  & 9 (same)                      & 3 & 7 & 8 & free $t_f$ (\code{cScale}) \\
\bottomrule
\end{tabular}
\end{center}

\noindent
Controls are $[\vbeta;s]$; the ELFO min-time solver pins $s\equiv1$ and carries
only $\vbeta$ ($n_u=3$), so the manifest's empty \code{thrRow} switches the
throttle-dependent checks off --- correctly, since an all-burn arc has no
switching structure to test. Empty manifest fields always mean ``skip that
check,'' never ``pretend it passed.''

\subsection{Multipliers done right}
\label{sec:duals}

The single most important implementation rule: multipliers come from the raw
low-level vector \code{opti.lam\_g}, indexed by \emph{constraint-row ranges
recorded at build time} (the \code{creg} registry each solver attaches under
\code{returnModel}). The convenience accessor \code{opti.dual()} returns duals
in CasADi's \emph{canonicalized} constraint orientation, which for a
row-wise-canonicalizing constraint (every collocation defect here) produces
entry-wise sign corruption --- the root cause of a 10--60$^\circ$ ``primer
anomaly'' that cost this project nine days of misdiagnosis in July 2026. The
full post-mortem is
\file{earth\_elliptic\_to\_geo/process/LESSONS\_DUAL\_EXTRACTION.md}; the
minimal reproduction is
\file{results/dual\_anomaly/diag\_optidual\_minimal.m}.

One genuine ambiguity remains: CasADi/IPOPT's overall Lagrangian sign
convention, $\Lag = f \pm \vlam_g^{\!\top} g$. It is resolved empirically,
once, for the whole multiplier vector:
\begin{equation}
s^\ast \;=\; \arg\min_{s\in\{+1,-1\}}
\bigl\| \nabla_z f + s\,A^{\!\top}\vlam_g \bigr\|_\infty,
\qquad A = \partial g/\partial z .
\label{eq:signres}
\end{equation}
A correct certificate drives one of the two residuals to machine precision and
leaves the other $O(1)$, so the resolution is unambiguous in practice
(e.g.\ $1.5\times10^{-14}$ vs $6.6$ on the Earth 10\,N row).

%=====================================================================
\section{The checks}
\label{sec:checks}

Each subsection states the condition formally, gives the intuition, names the
code, and says how to read the report line. Thresholds are collected in
Table~\ref{tab:thresholds} (Section~\ref{sec:report}).

%---------------------------------------------------------------
\subsection{KKT stationarity: the master residual}
\label{sec:kkt}

\begin{condition}[Lagrangian stationarity]
At a local minimizer $z^\ast$ of \eqref{eq:nlp} with multipliers $\vlam_g$,
\begin{equation}
\nabla_z \Lag \;=\; \nabla_z f(z^\ast) + s^\ast A^{\!\top}\vlam_g \;=\; 0 .
\label{eq:kkt}
\end{equation}
\end{condition}

\paragraph{Formal content.}
Equation~\eqref{eq:kkt} is the discrete image of \emph{several} continuous PMP
conditions at once. Stationarity with respect to an interior state column
$x_k$ yields
\begin{equation}
\frac{\Lambda_{k-1}-\Lambda_k}{w_k}
\;=\; A_k^{\!\top}
\underbrace{\frac{h_{k-1}\Lambda_{k-1}+h_k\Lambda_k}{h_{k-1}+h_k}}_{\textstyle \vlam(\sigma_k)},
\qquad w_k=\tfrac{h_{k-1}+h_k}{2},
\label{eq:adjoint}
\end{equation}
which is exactly the trapezoidal transcription of
$\dot\vlam=-\partial\Ham/\partial\vx$ --- including, in the CR3BP campaign, the
explicit time-dependence terms $\partial\Ham/\partial t$ that make the problem
non-autonomous --- and which is where the nodal map of
\S\ref{sec:costates} comes from. Stationarity with respect to a free horizon
variable ($\Delta L$, \code{cScale}) is the discrete condition attached to the
free span. So a machine-small $\|\nabla_z\Lag\|_\infty$ certifies the costate
dynamics and the horizon condition simultaneously, at collocation order,
without either being coded explicitly.

Two qualifications, from external review: \eqref{eq:adjoint} is the clean
adjoint recursion only at nodes where no \emph{other} constraint acts (at
boundary or path-constrained nodes the row also carries those multipliers);
and the free-span row is a discrete NLP condition, which coincides with the
physical $\Ham(t_f)$ transversality only once the reparameterization is
carried through --- see \S\ref{sec:timecostate}.

\paragraph{Is verifying KKT stationarity the same as verifying the costate
dynamics? Discretely yes --- and that is weaker than it sounds.}
\eqref{eq:adjoint} is an identity, not an approximation: the state-block
stationarity rows \emph{are} the discrete adjoint equations. But the costate
is \emph{defined} to be the multiplier, so the multipliers satisfy the
recursion by construction; we are not independently computing a costate and
testing it against an ODE. What the residual genuinely certifies is (i) that
the solver reached a KKT point rather than stopping early, and (ii) that the
multipliers in hand belong to the model we believe was solved --- which is
exactly what exposed the \code{opti.dual} corruption
(\S\ref{sec:duals}). It is silent on whether the discrete adjoint
\emph{approximates the continuous} one; that is a consistency/order question
governed by the mesh, invisible to the residual.

\paragraph{A caution about reading forward-integration failures.} The
analogous primal gap has been measured: re-integrating the exact direct
control from $\vx_0$ over $\sim$40 revs diverges by $\|r\|\sim3$,
$\|v\|\sim5$ while the defects read $10^{-14}$. This must \emph{not} be read
as ``the collocation solution is far from the continuous one.'' It is a
statement about the conditioning of an \emph{initial}-value problem over 40
revs --- the same instability that defeats indirect single shooting on this
problem --- whereas collocation solves a \emph{boundary}-value problem
stabilized by its terminal condition. Closeness in the BVP sense does not imply
closeness under forward propagation from a common initial state; that
distinction is the reason multiple shooting exists. The affirmative evidence
that the discrete solutions do approximate the continuous optimum is
cross-formulation agreement: the MEE/$L$-domain and Cartesian/Sundman stacks
share no independent variable, state vector, mesh or regularization, yet agree
at 10\,N on final mass to $\sim$0.03\% (1377.10 vs 1376.74\,kg), both inside
the published benchmark band \cite{HMG2004}. What genuinely remains uncertain
is finer: individual switch \emph{times} at deep rungs (0.2\,N runs
$\sim$3 nodes per switch, which is why switch counts are reported as
mesh bands), and the costates, for which no mesh-convergence study has ever
been run.

\paragraph{Intuition.}
``No first-order descent direction exists that respects the constraints.'' If
you could improve the objective by perturbing any state, control, or span
variable while staying feasible to first order, some component of
$\nabla_z\Lag$ would be nonzero. This is the strongest and least interpretable
check: it says \emph{extremal}, loudly, but points at nothing physical when it
fails.

\paragraph{Code.}
\file{foc\_check.m}, first block: builds
$\nabla_z f$ and $A$ by CasADi AD from \code{opti.f}, \code{opti.g}, evaluates
at the solved primal--dual pair, resolves the sign by \eqref{eq:signres}, and
reports \code{rep.kktStatInf} $=\|\nabla_z\Lag\|_\infty$ and
\code{rep.sLag} $=s^\ast$.

\paragraph{Interpretation.}
Healthy values are $10^{-14}$--$10^{-10}$; the observed range across all
certified rows is $1.5\times10^{-14}$ to $6.9\times10^{-10}$. The gate is a
deliberately loose $10^{-6}$: anything between $10^{-8}$ and $10^{-6}$
suggests a sloppily converged dual (re-run the warm re-solve tighter);
anything above $10^{-6}$ means the multipliers are \emph{not} a certificate
and every downstream line of the report is meaningless --- fix this first.

%---------------------------------------------------------------
\subsection{The minimum condition, direction part: tangential
$\partial\Lag/\partial\vbeta$}
\label{sec:direction}

\begin{condition}[Pointwise minimization over the thrust sphere]
At each node with $s_k>0$, $\vbeta_k$ minimizes $\Ham$ over
$\|\vbeta\|=1$, i.e.\ $\vbeta_k = -\,p_k/\|p_k\|$ where
$p_k = B(x_k)^{\!\top}\vlam_{v,k}$ is the primer vector \cite{Lawden1963}.
Equivalently, the \emph{tangential} component of the Hamiltonian gradient
vanishes:
\begin{equation}
\bigl(I-\vbeta_k\vbeta_k^{\!\top}\bigr)\,
\frac{\partial \Lag}{\partial \vbeta_k} \;=\; 0 .
\label{eq:tangential}
\end{equation}
\end{condition}

\paragraph{Formal content.}
$\Ham$ is linear in $\vbeta$ at fixed throttle, so minimizing over the unit
sphere aligns $\vbeta$ anti-parallel to the primer. The sphere constraint
$\|\vbeta\|^2=1$ carries its own multiplier, which is purely \emph{radial}
(its gradient is $2\vbeta$); projecting it out with $(I-\vbeta\vbeta^{\top})$
leaves a quantity that must vanish independently of that multiplier's value.

\paragraph{Intuition.}
``The thrust points exactly where the costates say it should.'' Any tangential
residual means rotating the thrust vector slightly would lower the Hamiltonian
--- a first-order improvement the optimizer left on the table.

\paragraph{Code.}
\file{foc\_check.m} direction block: extracts the $\vbeta$-rows of the full
Lagrangian gradient at each burn node, projects out the radial part, reports
\code{rep.dirTanMax} / \code{rep.dirTanMed} over burn nodes.

\paragraph{Interpretation --- with an honest caveat.}
Values at machine precision ($10^{-17}$--$10^{-16}$ observed) are expected.
Two limitations, recorded as pre-promotion items in the register's \S A6:

(i) \emph{Not independent.} The residual is bounded by
$\sqrt{n_\beta}\cdot$\code{kktStatInf} (here $n_\beta=3$, so $\sqrt3$) --- it
is a readable projection of check~\ref{sec:kkt}, not a separate test.
\textbf{External review correction (2026-07-25):} merely excluding the cone
multiplier before projecting does \emph{not} fix this, because the projector
annihilates the radial term anyway,
$P_{\vbeta}(q+2\mu\vbeta)=P_{\vbeta}q$ with $P_{\vbeta}=I-\vbeta\vbeta^{\!\top}$.
Genuine independence requires building $q_k$ from a \emph{different} gradient
--- e.g.\ defect duals only, or a separately reconstructed Hamiltonian --- and
comparing that against the solved direction.

(ii) \emph{Sign-blind.} $\vbeta=+p/\|p\|$ (the Hamiltonian \emph{maximizer})
satisfies \eqref{eq:tangential} identically. The branch is decided by the cone
multiplier: writing $\Lag=\Lag_0+\mu(\vbeta^{\!\top}\vbeta-1)$ with
$q+2\mu\vbeta=0$, the minimizer has $\vbeta^{\!\top}q<0$, equivalently
$\mu>0$; the maximizer has $\mu<0$ (reverse if the registered row is
$1-\vbeta^{\!\top}\vbeta$). Reporting the normalized anti-alignment
$-\vbeta^{\!\top}q/\|q\|$ is more informative than tangency alone.

Until both are addressed, minimality of the direction rests on the physical
layer (primer-alignment angle in \file{verify\_pmp\_mee}, the solver's own
\code{primerAlignDeg}) and on the sign law of \S\ref{sec:signlaw}, not on this
line.

%---------------------------------------------------------------
\subsection{The minimum condition, throttle part: the switching sign law}
\label{sec:signlaw}

\begin{condition}[Bang-bang law]
Write $\Ham = \Ham_0 + s\,S(t)$ (affine in the throttle at fixed direction;
$S$ is the \emph{switching function}). Along a fuel-optimal trajectory,
\begin{equation}
s^\ast(t) \;=\;
\begin{cases}
1, & S(t) < 0,\\
0, & S(t) > 0,
\end{cases}
\qquad\text{singular arcs where } S\equiv 0 .
\label{eq:banglaw}
\end{equation}
\end{condition}

\paragraph{Formal content.}
Because $\ell(s)=s$ is linear, the throttle minimizer sits at a \emph{bound},
not at a stationary point --- testing $\partial\Ham/\partial s = 0$ would be
the wrong condition. The discrete surrogate: form the Lagrangian gradient with
the throttle-bound multiplier rows \emph{removed},
\begin{equation}
S^d_k \;=\;
\Bigl[\nabla_z f + s^\ast A_{\setminus\mathrm{box}}^{\!\top}
\vlam_{g,\setminus\mathrm{box}}\Bigr]_{s_k\text{-row}},
\label{eq:Sd}
\end{equation}
which equals the continuous $S(\sigma_k)$ times a \emph{positive but
node-dependent} factor --- the local trapezoid weight $w_k$ times the clock
factor. Positivity means the \emph{sign} test in \eqref{eq:banglaw}, hence
\code{signPct}, is exact. It does \emph{not} mean $S^d$ is a faithful sampling
of $S$: because $w_k$ varies node to node, any statistic built on
\emph{differences} of $S^d$ across the mesh inherits the weighting. That is
the mechanism behind the regularity caveat in \S\ref{sec:sdot}, and it is why
$S^d$ must be deweighted ($\widehat S_k = S^d_k/w_k$) before differencing.
Removing the bound rows matters: at an active bound, KKT makes the \emph{full}
gradient row zero by absorbing $S^d_k$ into the bound multiplier --- the sign
information lives precisely in what the bound multiplier had to absorb.

\paragraph{Intuition.}
$S$ is the marginal Hamiltonian cost of thrust: ``would one more Newton-second
here help or hurt?'' Negative $\Rightarrow$ burn at full throttle; positive
$\Rightarrow$ coast. The check asks whether the solver's actual on/off pattern
agrees with its own multipliers' answer at \emph{every} node.

\paragraph{Code.}
\file{foc\_check.m} sign-law block: \code{rep.Sd} (the per-node values),
\code{rep.signPct} $=100\cdot\#\{k: (S^d_k<0) \Leftrightarrow s_k>0.5\}/(N{+}1)$.
Skipped (\code{'-{}-'}) when the manifest has no throttle row (min-time).

\paragraph{Interpretation.}
Gate: $\ge 99\%$. Certified rows read $100\%$. A few disagreeing nodes at
exactly the switch times are tolerable mesh artifacts (a node can straddle a
crossing); broad disagreement (the pre-fix era read 78\%) means the duals are
corrupted or the solution is not extremal.

\paragraph{Validity conditions (external review, 2026-07-25).}
Both external reviewers confirmed the recovery \eqref{eq:Sd} is exact --- the
active bound multiplier cancels $S^d_k$ identically, so removing its row
recovers magnitude \emph{and} sign --- but only under stated conditions:
the throttle cost must be \emph{affine} in $s$; the clock factor positive;
the removed rows the \emph{complete} set of constraints acting directly on
$s_k$; and the bound neither weakly active nor represented as coincident
lower/upper bounds (where the dual split is non-unique). At $\varepsilon>0$
the cost is quadratic in $s$ and $S^d$ is the derivative of the smoothed
objective, not an affine switching function: interior $s$ can legitimately
give $S^d=0$ and the sign carries no bang-bang implication.

\textbf{Fixed 2026-07-25.} \file{foc\_check.m} now resolves a throttle-cost
kind --- \code{opts.eps} (affine iff $\varepsilon=0$) overrides
\code{opts.throttleCostKind}, which overrides the manifest's own field --- and
\emph{skips} the whole bang-bang family (\code{signPct}, singular-arc count,
\code{sdotMinRel}) when the cost is not affine, reporting \code{'-{}-'} for
each and printing a NOTE. Every term in the advisory verdict is NaN-tolerant,
so a skipped check neither passes nor fails. Before this, the ELFO energy seed
($\varepsilon=1$) reported ``sign law 100\%'' and an ADVISORY PASS resting
partly on a meaningless number; it now passes on \code{kktStat} and
transversality alone, which are the conditions that actually apply at
$\varepsilon=1$.

%---------------------------------------------------------------
\subsection{Costate extraction: interval duals to nodal costates}
\label{sec:costates}

\begin{condition}[Discrete costate consistency]
The defect multipliers $\Lambda_k\in\mathbb R^{n_x}$ of \eqref{eq:defect} are
the discrete costates of interval $k$; the continuous costate sampled at an
interior node is their step-weighted average,
\begin{equation}
\vlam(\sigma_k) \;=\;
\frac{h_{k-1}\Lambda_{k-1} + h_k \Lambda_k}{h_{k-1}+h_k},
\qquad
\vlam(\sigma_1)=\Lambda_1,\quad \vlam(\sigma_{N+1})=\Lambda_N .
\label{eq:dualmap}
\end{equation}
\end{condition}

\paragraph{Formal content and intuition.}
A trapezoid defect couples nodes $k$ and $k{+}1$, so its multiplier ``lives on
the interval,'' not at a node --- the same way a finite-volume flux lives on a
face. The weighted average \eqref{eq:dualmap} is precisely the combination the
discrete adjoint recursion \eqref{eq:adjoint} pairs with the node-$k$ Jacobian
(equivalently, it is what appears in node-$k$ \emph{control} stationarity);
on a uniform mesh it reduces to the plain average of the two neighbors.
Derivation adjudicated in \file{process/DESIGN\_dual\_map.md}.

\paragraph{A distinction worth keeping (external review, 2026-07-25).}
The reviewers split here, and both were right about different questions.
\eqref{eq:dualmap} is the correct \emph{stationarity combination}. It is
\emph{not} the same thing as a linear \emph{interpolation} of interval-centred
samples to the node, which would carry the crossed weights
$(h_k\Lambda_{k-1}+h_{k-1}\Lambda_k)/(h_{k-1}+h_k)$. The two coincide on a
uniform mesh and differ on a refined one. This matters because
\code{rep.lam} is consumed \emph{both} ways: as a stationarity object (correct
as implemented) and as sampled function values feeding
\code{lamTimeCoV}, \code{lamTimeEnd}, and \code{lamMassEndRel} (where the
interpolation reading, and the endpoint caveat below, apply). Note also that no
exact nodal costate is recoverable from $\{\Lambda_k,h_k\}$ alone --- the exact
discrete covector involves the dynamics Jacobian, cf.\ \eqref{eq:adjoint}.

\paragraph{Code.}
\file{foc\_dual\_to\_costate.m} (generic in $n_x$; unit-tested at $n_x=9$ on a
non-uniform mesh). \file{foc\_check.m} applies it with the resolved sign:
\code{rep.lam} $= [\vlam(\sigma_1),\dots,\vlam(\sigma_{N+1})]$. This array is
an \emph{output for downstream consumers} (primer plots, PSR, the physical
verifiers), not itself a gate.

%---------------------------------------------------------------
\subsection{Transversality: free final mass}
\label{sec:transversality}

\begin{condition}[Free-terminal-state transversality]
The final mass is unconstrained and carries no terminal cost, so
\begin{equation}
\lambda_m(t_f) \;=\; 0 .
\label{eq:transv}
\end{equation}
\end{condition}

\paragraph{Intuition.}
$\lambda_m$ prices mass: ``what is one more kilogram at time $t$ worth to the
objective?'' At the end of the mission, leftover propellant buys nothing ---
its shadow price must be zero. A nonzero $\lambda_m(t_f)$ means the solution
is still trading mass against the objective at the final instant, i.e.\ it is
not done optimizing.

\paragraph{Code.}
\code{rep.lamMassEndRel} $= |\lambda_m(\sigma_{N+1})| / \max_k
|\lambda_m(\sigma_k)|$ --- \emph{relative}, because an absolute gate is
scale-dependent across campaigns. Skipped when the manifest declares the mass
not free at $t_f$.

\paragraph{Interpretation --- with an honest caveat.}
Gate: $\le 10^{-3}$. Most rows read $10^{-9}$--$10^{-4}$. Three rows (Earth
5\,N at $3.2\times10^{-3}$, 2.5\,N at $1.3\times10^{-3}$, CR3BP 5\,N at
$3.1\times10^{-3}$) miss the gate by $\sim3\times$ while their KKT residuals
are machine-tight. Two candidate readings, both recorded in the register:
(i) those are exactly the rows where warm re-solves find improved optima ---
under-converged mass costates on under-optimized rungs; or (ii) a
discretization artifact: by \eqref{eq:dualmap} the endpoint costate is the
\emph{one-sided} interval value, $\lambda_m$ at $\approx t_f - h/2$, which
carries an $O(h)\,|\dot\lambda_m|$ bias whenever the final arc burns.

\textbf{Both external reviewers independently confirmed reading (ii) is real},
and converged on the same cheap fix --- extrapolate from the last two interval
duals to the true endpoint,
\begin{equation}
\widehat\lambda(t_f)\;=\;\Lambda_N+\frac{h_N}{h_{N-1}+h_N}
\bigl(\Lambda_N-\Lambda_{N-1}\bigr),
\label{eq:endpointfix}
\end{equation}
with the symmetric formula at $\sigma_1$. \textbf{Implemented 2026-07-25}
(\file{foc\_dual\_to\_costate.m} second output \code{lamX}): the gate now uses
$\widehat\lambda(t_f)$ and \emph{reports the one-sided value beside it}, so a
marginal miss can be attributed rather than merely moved.

\paragraph{What that attribution found.} The two competing explanations
separated cleanly:
\begin{center}
\small
\begin{tabular}{@{}llll@{}}
\toprule
row & one-sided & endpoint-corrected & reading \\
\midrule
Earth 5\,N   & $3.189\times10^{-3}$ & $\mathbf{5.215\times10^{-4}}$ & endpoint bias; now PASS \\
CR3BP 5\,N   & $3.108\times10^{-3}$ & $\mathbf{6.437\times10^{-4}}$ & endpoint bias; now PASS \\
Earth 2.5\,N & $1.265\times10^{-3}$ & $1.265\times10^{-3}$ & \textbf{unchanged --- genuinely real} \\
\bottomrule
\end{tabular}
\end{center}
\noindent
The 2.5\,N row not moving \emph{at all} is the confirmation, not an anomaly:
the bias exists only where $\lambda_m$ is still changing at $t_f$, i.e.\ where
the final arc \emph{burns}. A row whose final arc coasts has
$\dot\lambda_m\approx0$ there and is untouched by \eqref{eq:endpointfix} --- so
its miss is a property of the solution, not of the discretization. That points
back at explanation (i): 2.5\,N is one of the rows a warm re-solve improves.

The more principled alternative remains open: form the exact discrete endpoint
covector $p_{N+1}=(I-\tfrac{h_N}{2}D_{N+1})^{\!\top}\Lambda_N$ plus terminal
terms, whose mass component is zero by terminal-node stationarity when the
final mass is genuinely free --- a discrete identity rather than an estimate.

%---------------------------------------------------------------
\subsection{The time costate: Hamiltonian conditions in dual form}
\label{sec:timecostate}

\begin{condition}[Hamiltonian constancy / terminal value]
With time carried as a state (row $t$ of $X$), its costate satisfies
$\dot\lambda_t = -\,\partial\Ham/\partial t$ and $\Ham \equiv -\lambda_t$
along the trajectory. Hence:
\begin{enumerate}[label=(\roman*),nosep]
\item autonomous dynamics, fixed $t_f$: $\lambda_t$ is \emph{constant}
      (generally nonzero) --- $\Ham$ is a first integral;
\item free $t_f$, objective $J=t(\tau_f)$ (min-time): transversality gives
      $\lambda_t(t_f) = \pm 1$ (sign fixed by the Lagrangian convention), and
      constancy then propagates that value backward;
\item non-autonomous (CR3BP: the Moon's phase is an explicit function of
      $t$): $\lambda_t$ is \emph{not} constant, and no constancy gate applies.
\end{enumerate}
\end{condition}

\paragraph{Intuition.}
$-\lambda_t$ \emph{is} the Hamiltonian, by the time-as-state trick: asking
``is $\Ham$ conserved?'' becomes ``is one costate row flat?'' --- a far easier
quantity to read off the NLP than a reconstructed $\Ham(t)$. The min-time
value condition is the classic ``the Hamiltonian equals minus the marginal
value of deadline time'': for $J=t_f$, one more second costs exactly one unit.

\paragraph{What this implementation does and does not establish
(external review, 2026-07-25).}
The identities $\Ham=-\lambda_t$ and $\lambda_t(t_f)=\pm1$ are classical for a
plain time-augmented system. They are \emph{not} derived here for a
Sundman-reparameterized system carrying a \code{cScale} free-time slack: that
would require fixing an extended Hamiltonian and costate convention and
propagating the reparameterization through them, which \file{foc\_check.m}
itself flags as pending (its \code{freetf-cscale} \code{horizonNote}). What
the code computes is the mapped defect-dual row and its statistics --- it never
forms $\Ham$, never verifies the identity, and never gates the endpoint value.
The ELFO reading $\lambda_t(t_f)=-1.000$ is therefore \emph{strong empirical
evidence} that the identity carries through the reparameterization, not a
verification of it. Completing that derivation is pre-promotion work.

\paragraph{Code.}
\code{rep.lamTimeCoV} (coefficient of variation of $\lambda_t$ over the nodes)
and \code{rep.lamTimeEnd}. This line is \emph{informational for every
campaign} --- it carries no PASS/FAIL --- and \code{rep.horizonNote} states
which case applies.

\paragraph{Interpretation.}
Earth fixed-$t_f$ rows show CoV $\sim10^{-9}$: textbook constancy. The ELFO
min-time anchor reads $\lambda_t(t_f) = -1.000$ --- the exact theoretical
free-$t_f$ transversality value, a strong independent confirmation obtained
with no min-time--specific code. Its CoV of $5.7\times10^{-2}$, however, is
\emph{not} small for an autonomous model; this is an open observation
(candidate explanations: the two-primary Sundman clock + \code{cScale}
structure distributes the horizon condition differently than the plain
reading assumes, or the anchor's duals are less converged than its primal) ---
see the register before leaning on that number.

%---------------------------------------------------------------
\subsection{Singular-arc detection}
\label{sec:singular}

\begin{condition}[No unexamined singular arcs]
On any interval where $S\equiv0$, the bang-bang law \eqref{eq:banglaw} is
silent and higher-order conditions (generalized Legendre--Clebsch) govern.
The check: no run of $\ge 3$ consecutive nodes with
$|S^d_k| < \max(10^{-6}\cdot\mathrm{scale},\,10^{-14})$,
where scale is the median $|S^d|$ over coast nodes.
\end{condition}

\paragraph{Intuition.}
A singular arc is a stretch where the marginal value of thrust is exactly
zero --- the optimizer is indifferent, and intermediate throttle can be
optimal. None are expected in these problems (and none have been seen:
\code{singularArcNodes = 0} on every row), but a solver stuck on a flat of
$S\approx0$ would silently produce garbage switch structure; this line makes
that loud instead.

\paragraph{Code/interpretation.}
\code{rep.singularArcNodes}; gate is $=0$. A nonzero count means: do not trust
the switch count or switch times of this row until the arc is examined with
the higher-order conditions.

%---------------------------------------------------------------
\subsection{Regular switching: $\dot S\neq 0$ at every switch}
\label{sec:sdot}

\begin{condition}[Transversal crossings]
At each throttle switch time $t_i$, the switching function crosses zero
transversally: $\dot S(t_i)\neq 0$. This is ingredient (ii) of the
Maurer--Osmolovskii sufficient conditions for a bang-bang local minimum
\cite{OsmolovskiiMaurer2012}, and it is what makes switch times well-defined,
isolated, and differentiable objects.
\end{condition}

\paragraph{Intuition.}
A switch where $S$ merely \emph{grazes} zero is fragile: an arbitrarily small
model perturbation can create, destroy, or slide it, and the finite-dimensional
switching-time problem that any second-order certificate is built on loses its
footing. Regular crossings are what let you speak of ``the 25 switches'' as a
stable, countable structure.

\paragraph{Code.}
\code{rep.sdotMinRel} $=\min_i |S^d_{k_i+1}-S^d_{k_i}| / \mathrm{scale}$ over
the switch nodes $k_i$; \code{rep.nSwitches}. Gate: $>10^{-3}$.

\paragraph{Interpretation --- with the register's caveat.}
\textbf{Fixed 2026-07-25 --- and the original readings were artifacts.} The
un-normalized statistic reported tulip flagship $1.4\times10^{-7}$ (25
switches) and ELFO $1.33\times$ front row $4.5\times10^{-5}$ (50 switches),
which looked like confirmation of the node-grazing-switch behavior the tulip
campaign already suspected. Both external reviewers independently judged those
numbers \emph{unable to diagnose grazing}, because $S^d_k$ carries the local
trapezoid weight $w_k$ and the difference was not divided by the step: on
refined or PSR-split meshes (which densify \emph{exactly at switches}) the
statistic shrinks with $h$ even for a perfectly regular crossing. They were
right --- under the corrected statistic the same two rows read
$\mathbf{2.70\times10^{1}}$ and $\mathbf{1.19\times10^{2}}$, comfortably
transversal, and both now PASS. \emph{The grazing reading was entirely a mesh
artifact.}

A dedicated regression test (\file{tests/test\_foc\_mesh\_invariance.m}) locks
this in both directions on the toy problem: across a $4\times$ refinement the
normalized statistic holds ($2.727\to2.775$, ratio $1.02$) while the legacy one
falls $4.82\times10^{-2}\to1.18\times10^{-2}$ --- a ratio of $4.1$, i.e.\ it was
reporting $h$.

The replacement --- deweight, then take a true divided difference, then
non-dimensionalize:
\begin{equation}
\widehat S_k=\frac{S^d_k}{w_k},\qquad
D_i=\left|\frac{\widehat S_{k_i+1}-\widehat S_{k_i}}
{\sigma_{k_i+1}-\sigma_{k_i}}\right|,\qquad
R_i=\frac{(\sigma_{N+1}-\sigma_1)\,D_i}{S_{\mathrm{ref}}},
\label{eq:sdotfix}
\end{equation}
with $w_1=h_1/2$, $w_k=(h_{k-1}+h_k)/2$, $w_{N+1}=h_N/2$, and
$S_{\mathrm{ref}}$ the median $|\widehat S_j|$ over coast nodes outside a
one-node neighborhood of any switch. Report $\min_i R_i$: dimensionless and
invariant to local refinement for a smooth transversal crossing. The same
deweighting is applied to the singular-arc test of \S\ref{sec:singular}, which
otherwise reports false positives in densified regions.

\textbf{Caveat that survives the fix.} One reviewer adds --- correctly --- that
no threshold on a \emph{single} mesh is defensible: $R_i$ being large on one
mesh does not establish transversality. Regularity should be asserted only if
$R_i$ stays above the floor on \emph{two} meshes, and the $10^{-3}$ gate is
provisional until recalibrated on deweighted data. So the current PASSes on
this line mean ``no evidence of grazing,'' not ``proven regular'' --- which
matters, because this is Maurer ingredient~(ii) and a second-order certificate
will lean on it.

%---------------------------------------------------------------
\subsection{The second-order line: IPOPT native inertia ($\delta_w$)}
\label{sec:inertia}

Strictly beyond first order, but printed in the same report because it is
free. IPOPT's inertia-controlled linear solver adds a Hessian regularization
$\delta_w>0$ \emph{only} when the reduced Hessian is indefinite at the current
iterate. If the solver converges with $\delta_w=0$ over its final iterations,
no inertia correction was needed at those iterates. The interpreter is
\file{foc\_ipopt\_inertia.m}, reading the $\delta_w$ history that every
production solver now records (\code{out.regHistory}); verdicts are
\code{LOCAL MIN} / \code{NOT CERTIFIED} / \code{NO-DATA}.

\paragraph{How much this is worth (external review, 2026-07-25).}
Read it as an \emph{inertia-regularization observation}, not a second-order
certificate. $\delta_w=0$ over a tail says the correction was not applied at
those \emph{barrier} iterates; it does not by itself prove positive
definiteness of the final reduced Hessian on the critical cone, account for
barrier/active-set limiting behavior, or establish strictness. Even read
generously the ceiling is low: at $\varepsilon>0$ the problem is strictly
convex in throttle (the easy case), and at $\varepsilon=0$ the barrier supplies
the throttle-direction curvature, so the most it can support is a \emph{weak},
non-strict local minimum --- strictness lives in the switching times, which is
exactly the still-open Maurer curvature test. The \code{certLocalMin} /
\code{LOCAL MIN} naming (inherited from the tulip campaign's
\file{psr\_ipopt\_certify.m}) overstates this and should be relabeled
campaign-wide.

%=====================================================================
\section{Running the layer and reading the report}
\label{sec:report}

\subsection{Entry points per campaign}

All commands assume MATLAB R2025b, run from the named folder, after
\code{setup\_paths} and
\code{addpath(fullfile(getenv('HOME'),'casadi-3.7.0'))}. Each entry point
loads a certified artifact, \emph{warm re-solves at its saved primal} (to
attach the live CasADi model and re-derive multipliers under the corrected
extraction), guards that re-solve on the certified quantities (solver success,
defect $\le10^{-8}$, one-sided final mass --- or non-degraded $t_f$ for
min-time), then runs \code{foc\_check} and prints \code{foc\_report}.

\begin{itemize}[leftmargin=2em]
\item \textbf{Earth two-body} (\file{earth\_elliptic\_to\_geo/direct/}):\\
\code{[rep, ver] = run\_foc\_mee('results/MEE\_M2\_10N.mat');}\\
Ladder-wide sweep (all nine certified rows, generic + physical gates side by
side): \code{T = run\_verify\_pmp\_all();}
\item \textbf{Earth CR3BP} (\file{earth\_elliptic\_to\_geo\_CR3BP/direct/}):\\
\code{ver = verify\_cr3bp\_pmp(struct('thrustN',10,'phi0',0));}
(FOC block runs automatically; result in \code{ver.foc}).
\item \textbf{GTO$\to$tulip} (\file{GTO\_tulip/direct/sundman\_minfuel/}):\\
\code{rep = run\_foc\_tulip();} --- defaults to the certified 25-switch
flagship; also runs the independent least-squares costate reconstruction
(\code{certify\_minfuel\_pmp}) and records both in \code{rep.crossCheck}.
\item \textbf{GTO$\to$ELFO} (\file{GTO\_ELFO/direct/elfo/}):
\begin{small}
\begin{verbatim}
rep = run_foc_elfo('mintime', fullfile('results','mintime_elfo.mat'));
rep = run_foc_elfo('fuel', ...
       fullfile('results','minfuel_ELFO_tf1p330_sw50_minEps0.mat'));
\end{verbatim}
\end{small}
\end{itemize}

\noindent
In addition, the production drivers (\code{run\_gergaud},
\code{run\_cr3bp\_geo}, \code{run\_psr}, \code{run\_elfo\_minfuel},
\code{gen\_elfo\_mintime}) emit the report automatically after a live
certifying solve, try/catch-guarded so a report failure can never fail a
solve. Each run also saves a sidecar
\code{results/foc\_<tag>.mat} containing the full \code{rep} struct for later
inspection.

\subsection{Anatomy of the report}

A real block, from the Earth 2.5\,N row (July 2026) --- chosen deliberately
because reading a \emph{failing} report correctly is the skill this section
needs to teach:

\begin{small}
\begin{verbatim}
========== FIRST-ORDER OPTIMALITY REPORT: MEE_M2_2p5N ==========
 KKT stationarity  ||grad_x L||_inf :  1.166e-11   (sign s=+1)   PASS
 Min condition (throttle sign law)  :    100.00 %                PASS
 Transversality |lam_m(tf)| rel     :  1.265e-03                 FAIL
 Singular-arc nodes                 :          0                 PASS
 Regular switching min|Sdot| rel    :  7.247e+00   (76 switches) PASS
   ^ transversality, legacy one-sided :  1.265e-03   [I5 endpoint fix]
   ^ regular switching, legacy raw    :  9.888e-03   [I1 mesh fix]
 ADVISORY verdict: FAIL   (report-only burn-in: does not alter certified status)
\end{verbatim}
\end{small}

\noindent
Line-by-line: the row is a machine-tight extremal (\S\ref{sec:kkt}) whose
thrust program agrees with its own multipliers at every node
(\S\ref{sec:signlaw}), has no singular arcs (\S\ref{sec:singular}), and whose
76 switches all cross transversally (\S\ref{sec:sdot}) --- but it has
\emph{not} priced its propellant to zero at $t_f$
(\S\ref{sec:transversality}). The advisory verdict is the conjunction of the
gated lines, so it is FAIL solely on that one; per the burn-in policy this
changes nothing about the row's certified status --- it files a finding.

\paragraph{Reading the two \texttt{\^{}} attribution lines.} These carry the
diagnosis. The regular-switching line moved $9.888\times10^{-3}\to7.247$ under
the I1 mesh normalization --- a large move, so the legacy number was a
discretization artifact and there was never a grazing problem here. The
transversality line did \emph{not} move ($1.265\times10^{-3}$ both ways), so
the I5 endpoint bias is \emph{not} the explanation and the miss is real. A
large gap means ``the old number was an artifact''; a zero gap means ``this
finding survives correction.'' That is what turns a FAIL from an alarm into a
diagnosis --- and it is why both values are printed rather than the corrected
one alone.

The direction line (\S\ref{sec:direction}), the time-costate line
(\S\ref{sec:timecostate}), the horizon note, and the $\delta_w$ line
(\S\ref{sec:inertia}) appear when applicable; a \code{'-{}-'} status always
means ``check skipped by the manifest or value not applicable,'' never a
silent pass.

\begin{table}[h]
\centering
\small
\begin{tabular}{@{}llll@{}}
\toprule
report line & \code{rep} field & gate & skipped when \\
\midrule
KKT stationarity & \code{kktStatInf} & $\le 10^{-6}$ & never \\
Min condition (direction) & \code{dirTanMax} & $\le 10^{-6}$ & \code{dirRows} empty \\
Min condition (sign law) & \code{signPct} & $\ge 99\,\%$ & \code{thrRow} empty (min-time) \\
Transversality (mass) & \code{lamMassEndRel} & $\le 10^{-3}$ & mass not free at $t_f$ \\
Time-costate CoV / end & \code{lamTimeCoV}, \code{lamTimeEnd} & informational & \code{timeRow} empty \\
Singular-arc nodes & \code{singularArcNodes} & $=0$ & \code{thrRow} empty \\
Regular switching & \code{sdotMinRel} & $> 10^{-3}$ & \code{thrRow} empty \\
IPOPT inertia $\delta_w$ & \code{ipopt.certLocalMin} & informational & no \code{regHistory} \\
\midrule
\multicolumn{4}{@{}l@{}}{\emph{attribution companions (printed beside the gated value)}}\\
transversality, one-sided & \code{lamMassEndRelOneSided} & --- & as above \\
regular switching, raw & \code{sdotMinRelLegacy} & --- & as above \\
deweighted switching fn & \code{Sdeweighted} & --- & \code{thrRow} empty \\
throttle-cost kind & \code{throttleCostKind}, \code{bangBangChecksRun} & --- & never \\
\bottomrule
\end{tabular}
\caption{Report lines, their \code{rep} fields, and default thresholds
(\file{foc\_check.m} defaults; \file{foc\_report.m} prints matching labels).
The advisory verdict \code{rep.pass} is the conjunction of the gated rows.}
\label{tab:thresholds}
\end{table}

\subsection{Policy and provenance}

The layer runs in \textbf{report-only burn-in}: \code{rep.pass} is advisory,
every printed block carries the disclaimer, and no certified row can be
demoted by a FAIL. Promotion to a hard gate is a recorded future decision,
blocked on the \S A6 checklist of
\file{orbit\_transfer/OPTIMALITY\_CERTIFICATION.md} (mesh-normalize
$\dot S$; make the direction check independent and sign-aware; resolve the
transversality endpoint bias). That register --- not this note --- is the
living source of per-row results and open findings; this note explains what
the instruments measure and how to hold them.

\begin{thebibliography}{9}

\bibitem{HMG2004}
T.~Haberkorn, P.~Martinon, and J.~Gergaud,
``Low-thrust minimum-fuel orbital transfer: a homotopic approach,''
\emph{Journal of Guidance, Control, and Dynamics}, 27(6), 2004.

\bibitem{BertrandEpenoy2002}
R.~Bertrand and R.~\'Ep\'enoy,
``New smoothing techniques for solving bang-bang optimal control problems ---
numerical results and statistical interpretation,''
\emph{Optimal Control Applications and Methods}, 23(4), 2002.

\bibitem{Lawden1963}
D.~F. Lawden,
\emph{Optimal Trajectories for Space Navigation},
Butterworths, London, 1963.

\bibitem{OsmolovskiiMaurer2012}
N.~P. Osmolovskii and H.~Maurer,
\emph{Applications to Regular and Bang-Bang Control: Second-Order Necessary
and Sufficient Optimality Conditions in Calculus of Variations and Optimal
Control}, SIAM, Philadelphia, 2012.

\bibitem{Betts2010}
J.~T. Betts,
\emph{Practical Methods for Optimal Control and Estimation Using Nonlinear
Programming}, 2nd ed., SIAM, Philadelphia, 2010.

\end{thebibliography}

\end{document}
